% =====================================================================
%  Tabulka symbolů
% =====================================================================
\section{Tabulka symbolů}\label{sec:symtable}

Na základě vybrané varianty vv-BVS jsme museli tabulku symbolů implementovat jako výškově vyvážený binární
vyhledávací strom (zvolili jsme AVL strom~\cite{ial06vyhledavani}). Klíčem do tabulky symbolů je jméno identifkátoru, uzly se řadí lexikograficky
pomocí \fce{strcmp}.

\subsection{Struktura}
Funkce a proměnné jsou uloženy v~tabulce symbolů odděleně. Funkce (uživatelské i~vestavěné) jsou v~jednom globálním
stromu, protože jsou viditelné pro celý program. Proměnné se ukládají do zásobníku rozsahů, který je implementován
jako jednosměrně vázaný seznam. Každý rozsah by měl odpovídat jednomu bloku zdrojového kódu
(má vlastní strom a hloubku na zásobníku). Sémantická analýza vždy při vstupu do bloku vloží rozsah (\fce{symtable\_push\_scope}) a když jej opouští, volá se
\fce{symtable\_pop\_scope}.

\subsection{Obsah symbolu}
Každý uzel obsahuje strukturu \texttt{symbol\_t}. U~proměnné obsahuje typ, příznak konstanty a jedinečný identifikátor. U~funkce symbol obsahuje navíc ještě návratový typ,
typy a jména parametrů, příznak vestavěné funkce, příznak \texttt{@main}, stav kontroly těla (kapitola~\ref{sec:semantika}) a příznak odvození
návratového typu a ukazatel na definici funkce v~AST.

\subsection{Operace se stromem}
Vkládaní probíhá rekurzivně sestupem. Když se z~rekurze vrací, přepočítá se výška a strom se případně vyváží na základě rozdílu výšek podstromů.
Vyvážení se provádí aplikací jedné ze čtyř možných rotací (LL, LR, RL, RR). Když uzel se stejným klíčem už existuje, funkce která vkládaní provádí,
vrátí \texttt{ST\_EXISTS}. V~sémantice se z toho stane chyba~4.

Vyhledávání probíhá následovně: proměnná se hledá od vrcholu zásobníku směrem dolů a
vratí se první nalezená hodnota - ta z nejvnitřnějšího bloku (funguje tak zastíňení). Pro kontrolu redefinice prohledá funkce \fce{symtable\_lookup\_local} jen aktuální rozsah.

Proměnné zanikají vždy s~celým blokem (při odstranění rozsahu se uvolňuje celý jeho strom).
Vkládaní i~vyhledávání má časovou složitost $O(\log n)$.

Součástí tabulky symbolů je taky čítač, který přiřazuje jedinečná čísla symbolům.
Vkládaní vrací ukazatel na vložený symbol, aby ho volající mohl hned doplnit.

